% ======================================================================
% Topic: Ai
% ======================================================================

\element{ai}{
  \begin{question}{Q001}
    ข้อใดอธิบาย Supervised Learning ได้ถูกต้องที่สุด
    \begin{choices}
        \wrongchoice{ปล่อยให้โมเดลค้นหาโครงสร้างข้อมูลเองโดยไม่มีคำตอบ}
        \wrongchoice{โมเดลเรียนรู้จากรางวัลและการลงโทษจากสภาพแวดล้อม}
        \correctchoice{โมเดลเรียนรู้จากตัวอย่างที่ถูกต้อง}
        \wrongchoice{โมเดลสุ่มทดลองโดยไม่ต้องประเมินผล}
    \end{choices}
  \end{question}
}

\element{ai}{
  \begin{question}{Q002}
    ลักษณะสำคัญของ Unsupervised Learning คืออะไร
    \begin{choices}
        \wrongchoice{โมเดลสุ่มทดลองโดยไม่ต้องประเมินผล}
        \wrongchoice{เรียนรู้จากรางวัลและโทษ}
        \correctchoice{ไม่มีคำตอบให้ และให้โมเดลค้นหาโครงสร้างข้อมูลเอง}
        \wrongchoice{ใช้กฎที่มนุษย์เขียนไว้ล้วน ๆ}
    \end{choices}
  \end{question}
}

% \element{ai}{
%   \begin{question}{Q003}
%     ผลกระทบต่อแรงงานจาก AI ควรถูกมองอย่างไร
%     \begin{choices}
%         \wrongchoice{เป็นผลบวกเสมอ}
%         \wrongchoice{เป็นผลลบเสมอ}
%         \correctchoice{อาจเป็นได้ทั้งบวกและลบ ขึ้นอยู่กับการจัดการ}
%         \wrongchoice{ไม่มีผลกระทบจริง}
%     \end{choices}
%   \end{question}
% }

% \element{ai}{
%   \begin{question}{Q004}
%     เหตุใดงาน WhiteCollar บางประเภทจึงมีความเสี่ยงสูงในการถูกแทนที่ด้วย AI
%     \begin{choices}
%         \wrongchoice{เพราะเป็นงานที่ใช้แรงงาน}
%         \correctchoice{เพราะสามารถทำงานซ้ำ ๆ ตามขั้นตอนที่กำหนดได้}
%         \wrongchoice{เพราะไม่ต้องใช้ทักษะการคิดวิเคราะห์}
%         \wrongchoice{เพราะไม่ต้องใช้คอมพิวเตอร์}
%     \end{choices}
%   \end{question}
% }

% \element{ai}{
%   \begin{question}{Q003}
%     คำว่า ``Summarize" จัดอยู่ในส่วนใดของ Prompt Structure
%     \begin{choices}
%         \correctchoice{Instruction}
%         \wrongchoice{Context}
%         \wrongchoice{Input}
%         \wrongchoice{Output}
%     \end{choices}
%   \end{question}
% }

% \element{ai}{
%   \begin{question}{Q006}
%     ข้อความ ``a presentation on recent advances in machine learning" จัดอยู่ในส่วนใด
%     \begin{choices}
%         \wrongchoice{Instruction}
%         \correctchoice{Context}
%         \wrongchoice{Input}
%         \wrongchoice{Output}
%     \end{choices}
%   \end{question}
% }

% \element{ai}{
%   \begin{question}{Q007}
%     ตัวอย่าง ``Translate Hello to Thai" จัดอยู่ในรูปแบบใด
%     \begin{choices}
%         \wrongchoice{Few Shot}
%         \correctchoice{Zero Shot}
%         \wrongchoice{Domain Expertise}
%         \wrongchoice{Contextual Guidance}
%     \end{choices}
%   \end{question}
% }

% \element{ai}{
%   \begin{question}{Q008}
%     ข้อใดเป็นตัวอย่างของ Few Shot Prompting
%     \begin{choices}
%         \wrongchoice{Translate ``Hello" to Thai}
%         \correctchoice{Hello: สวัสดีเจ้า / Good bye: ไปก่อนเน้อเจ้า / Goodnight:…}
%         \wrongchoice{Write a short essay on AI}
%         \wrongchoice{Summarize the article in 100 words}
%     \end{choices}
%   \end{question}
% }

% \element{ai}{
%   \begin{question}{Q009}
%     ตัวอย่างคำถาม ``How to get a good grade" ใน Tree of Thought ใช้เพื่ออะไร
%     \begin{choices}
%         \wrongchoice{สาธิตการคำนวณเชิงตัวเลข}
%         \correctchoice{สาธิตการวิเคราะห์แบบหลายมุมมอง}
%         \wrongchoice{สาธิตการใช้ ZeroShot Prompting}
%         \wrongchoice{สาธิตการลด Bias}
%     \end{choices}
%   \end{question}
% }

% \element{ai}{
%   \begin{question}{Q010}
%     Management for Sustainability มุ่งเน้นเรื่องใด
%     \begin{choices}
%         \wrongchoice{การจัดเก็บภาษีท้องถิ่น}
%         \correctchoice{การมีแผนและกลไกเพื่อความยั่งยืน และการมีส่วนร่วมของประชาชน}
%         \wrongchoice{การขยายเมืองออกไปสู่ชานเมือง}
%         \wrongchoice{การสร้างสนามกีฬาในเมือง}
%     \end{choices}
%   \end{question}
% }

% \element{ai}{
%   \begin{question}{Q013}
%     ผลกระทบต่อแรงงานจาก AI ควรถูกมองอย่างไร
%     \begin{choices}
%         \wrongchoice{เป็นผลบวกเสมอ}
%         \wrongchoice{เป็นผลลบเสมอ}
%         \wrongchoice{อาจเป็นได้ทั้งบวกและลบ ขึ้นอยู่กับการจัดการ}
%         \wrongchoice{ไม่มีผลกระทบจริง}
%     \end{choices}
%   \end{question}
% }

\element{ai}{
  \begin{question}{Q004}
    เหตุใดงานบางประเภทจึงมีความเสี่ยงสูงในการถูกแทนที่ด้วย AI
    \begin{choices}
        \wrongchoice{เพราะเป็นงานที่ใช้แรงงาน}
        \correctchoice{เพราะสามารถทำงานซ้ำ ๆ ตามขั้นตอนที่กำหนดได้}
        \wrongchoice{เพราะไม่ต้องใช้ทักษะการคิดวิเคราะห์}
        \wrongchoice{เพราะไม่ต้องใช้คอมพิวเตอร์}
    \end{choices}
  \end{question}
}

% \element{ai}{
%   \begin{question}{Q015}
%     คำว่า ``Summarize" จัดอยู่ในส่วนใดของ Prompt Structure
%     \begin{choices}
%         \wrongchoice{Instruction}
%         \wrongchoice{Context}
%         \wrongchoice{Input}
%         \wrongchoice{Output}
%     \end{choices}
%   \end{question}
% }

% \element{ai}{
%   \begin{question}{Q016}
%     ข้อความ ``a presentation on recent advances in machine learning" จัดอยู่ในส่วนใด
%     \begin{choices}
%         \wrongchoice{Instruction}
%         \wrongchoice{Context}
%         \wrongchoice{Input}
%         \wrongchoice{Output}
%     \end{choices}
%   \end{question}
% }

% \element{ai}{
%   \begin{question}{Q017}
%     ตัวอย่าง ``Translate ‘Hello’ to Thai" จัดอยู่ในรูปแบบใด
%     \begin{choices}
%         \wrongchoice{Few Shot}
%         \wrongchoice{Zero Shot}
%         \wrongchoice{Domain Expertise}
%         \wrongchoice{Contextual Guidance}
%     \end{choices}
%   \end{question}
% }

% \element{ai}{
%   \begin{question}{Q018}
%     ข้อใดเป็นตัวอย่างของ Few Shot Prompting
%     \begin{choices}
%         \wrongchoice{Translate ``Hello" to Thai}
%         \wrongchoice{Hello: สวัสดีเจ้า / Good bye: ไปก่อนเน้อเจ้า / Goodnight:…}
%         \wrongchoice{Write a short essay on AI}
%         \wrongchoice{Summarize the article in 100 words}
%     \end{choices}
%   \end{question}
% }

% \element{ai}{
%   \begin{question}{Q019}
%     ตัวอย่างคำถาม ``How to get a good grade" ใน Tree of Thought ใช้เพื่ออะไร
%     \begin{choices}
%         \wrongchoice{สาธิตการคำนวณเชิงตัวเลข}
%         \wrongchoice{สาธิตการวิเคราะห์แบบหลายมุมมอง}
%         \wrongchoice{สาธิตการใช้ ZeroShot Prompting}
%         \wrongchoice{สาธิตการลด Bias}
%     \end{choices}
%   \end{question}
% }

% \element{ai}{
%   \begin{question}{Q020}
%     Management for Sustainability มุ่งเน้นเรื่องใด
%     \begin{choices}
%         \wrongchoice{การจัดเก็บภาษีท้องถิ่น}
%         \wrongchoice{การมีแผนและกลไกเพื่อความยั่งยืน และการมีส่วนร่วมของประชาชน}
%         \wrongchoice{การขยายเมืองออกไปสู่ชานเมือง}
%         \wrongchoice{การสร้างสนามกีฬาในเมือง}
%     \end{choices}
%   \end{question}
% }

% \element{ai}{
%   \begin{question}{Q021}
%     Bruce Henderson (1975) เสนอแนวคิดเรื่องใด
%     \begin{choices}
%         \wrongchoice{Value Chain}
%         \correctchoice{Scale curves, Experience curves}
%         \wrongchoice{Five Forces Model}
%         \wrongchoice{Balanced Scorecard}
%     \end{choices}
%   \end{question}
% }

% \element{ai}{
%   \begin{question}{Q022}
%     Michael Porter (1980) เน้นแนวคิดใด
%     \begin{choices}
%         \correctchoice{Value Chain}
%         \wrongchoice{Balanced Scorecard}
%         \wrongchoice{Stakeholder Theory}
%         \wrongchoice{Game Theory}
%     \end{choices}
%   \end{question}
% }

% \element{ai}{
%   \begin{question}{Q023}
%     ข้อใดคือความแตกต่างหลักของแนวคิด Henderson (1975) และ Porter (1980)
%     \begin{choices}
%         \correctchoice{Henderson = เน้นปริมาณ/ประสบการณ์กับต้นทุน, 	Porter = เน้นห่วงโซ่คุณค่าและโครงสร้างธุรกิจ}
%         \wrongchoice{Henderson = เน้น Value Chain, Porter = เน้น Economies of Mass}
%         \wrongchoice{Henderson = เน้น Balanced Scorecard, Porter = เน้น Game Theory}
%         \wrongchoice{ไม่มีความต่าง}
%     \end{choices}
%   \end{question}
% }

\element{ai}{
  \begin{question}{Q005}
    ข้อใดคือความหมายของ Artificial Intelligence (AI)
    \begin{choices}
        \wrongchoice{ความสามารถของระบบที่เรียนรู้ได้เองโดยไม่ต้องเขียนโปรแกรมชัดเจน}
        \wrongchoice{การเรียนรู้ที่อาศัยโครงข่ายประสาทเทียมหลายชั้น}
        \correctchoice{โปรแกรมที่สามารถรับรู้ เหตุผล ลงมือทำ และปรับตัวได้}
        \wrongchoice{การวิเคราะห์ข้อมูลเพื่อจัดกลุ่ม (Clustering)}
    \end{choices}
  \end{question}
}

\element{ai}{
  \begin{question}{Q006}
    ลักษณะสำคัญของ Reinforcement Learning คืออะไร
    \begin{choices}
        \wrongchoice{โมเดลสุ่มทดลองโดยไม่ต้องประเมินผล}
        \correctchoice{เรียนรู้จากการให้รางวัลและการลงโทษ}
        \wrongchoice{จัดกลุ่มข้อมูลเองโดยไม่รู้คำตอบ}
        \wrongchoice{ใช้กฎที่เขียนไว้}
    \end{choices}
  \end{question}
}

\element{ai}{
  \begin{question}{Q007}
    หลักการสำคัญของ K-means clustering คืออะไร
    \begin{choices}
        \correctchoice{หาค่าเฉลี่ย ของแต่ละกลุ่มและจัดข้อมูลไปยังกลุ่มที่ใกล้ที่สุด}
        \wrongchoice{ใช้กฎที่มนุษย์เขียนเพื่อจัดกลุ่ม}
        \wrongchoice{ทำนายค่าตัวเลขต่อเนื่อง}
        \wrongchoice{ใช้รางวัลและโทษเพื่อฝึก AI}
    \end{choices}
  \end{question}
}

\element{ai}{
  \begin{question}{Q008}
    AI และคอมพิวเตอร์มีผลกระทบต่อการจ้างงานในลักษณะใดบ้าง
    \begin{choices}
        \wrongchoice{ใช้แทนแรงงานในบางอาชีพ}
        \wrongchoice{การเปลี่ยนแปลงจะเกิดขึ้นอย่างรวดเร็ว}
        \wrongchoice{ลดจำนวนแรงงาน}
        \lastchoices
        \correctchoice{ถูกทุกข้อ}
    \end{choices}
  \end{question}
}

\element{ai}{
  \begin{question}{Q009}
    ความแตกต่างที่ถูกต้องระหว่าง Automation กับ Augmentation คืออะไร
    \begin{choices}
        \correctchoice{Automation = แทนที่งานบางส่วนด้วย AI และ หุ่นยนต์, Augmentation = ใช้ AI ผู้ช่วยเพื่อเพิ่มประสิทธิภาพงาน}
        \wrongchoice{ทั้งสองหมายถึงการเลิกจ้างพนักงานทั้งหมด}
        \wrongchoice{Automation = เพิ่มจำนวนพนักงาน, Augmentation = ลดคุณภาพงาน}
        \wrongchoice{Automation = ใช้คนล้วน, Augmentation = ไม่ใช้คนเลย}
    \end{choices}
  \end{question}
}

% \element{ai}{
%   \begin{question}{Q010}
%     เหตุใดงาน WhiteCollar บางประเภทจึงมีความเสี่ยงสูงในการถูกแทนที่ด้วย AI
%     \begin{choices}
%         \wrongchoice{เพราะเป็นงานที่ใช้แรงงาน}
%         \correctchoice{เพราะสามารถทำงานซ้ำ ๆ ตามขั้นตอนที่กำหนดได้}
%         \wrongchoice{เพราะไม่ต้องใช้ทักษะการคิดวิเคราะห์}
%         \wrongchoice{เพราะไม่ต้องใช้คอมพิวเตอร์}
%     \end{choices}
%   \end{question}
% }

\element{ai}{
  \begin{question}{Q011}
    ข้อใดต่อไปนี้จัดอยู่ในกลุ่ม AI Services
    \begin{choices}
        \wrongchoice{Netflix}
        \correctchoice{Suno.com}
        \wrongchoice{Canva.com}
        \wrongchoice{MS Offices}
    \end{choices}
  \end{question}
}

\element{ai}{
  \begin{question}{Q012}
    ตัวอย่างของ Integrated AI คือข้อใด
    \begin{choices}
        \wrongchoice{Runwayml.com}
        \wrongchoice{Napkin.ai}
        \correctchoice{Notion.com}
        \wrongchoice{Lovable.dev}
    \end{choices}
  \end{question}
}

\element{ai}{
  \begin{question}{Q013}
    Agentic AI คืออะไร
    \begin{choices}
        \wrongchoice{ไม่สามารถใช้ฐานข้อมูลภายนอกได้}
        \correctchoice{ตัดสินใจเลือกเครื่องมือเพื่อแก้ปัญหาได้เอง}
        \wrongchoice{ใช้ได้เฉพาะงานแปลภาษา}
        \wrongchoice{ต้องเขียนโค้ดทุกขั้นตอน}
    \end{choices}
  \end{question}
}

% \element{ai}{
%   \begin{question}{Q014}
%     ``Translate this English sentence into French." เป็นตัวอย่างของเทคนิคใด
%     \begin{choices}
%         \wrongchoice{Contextual Guidance}
%         \correctchoice{Task Specific}
%         \wrongchoice{Domain Expertise}
%         \wrongchoice{Framing.}
%     \end{choices}
%   \end{question}
% }

% \element{ai}{
%   \begin{question}{Q015}
%     ``Explain the causes, symptoms, and treatments of hypothyroidism, including the latest research and medical guidelines." จัดอยู่ในเทคนิคใด
%     \begin{choices}
%         \correctchoice{Domain Expertise}
%         \wrongchoice{Task Specific}
%         \wrongchoice{Framing}
%         \wrongchoice{Contextual Guidance}
%     \end{choices}
%   \end{question}
% }

% \element{ai}{
%   \begin{question}{Q016}
%     ``Write a 100word paragraph on leadership traits without favoring any gender." เป็นตัวอย่างของเทคนิคใด
%     \begin{choices}
%         \wrongchoice{User Feedback Loop}
%         \correctchoice{Bias Mitigation}
%         \wrongchoice{Task Specific}
%         \wrongchoice{Contextual Guidance}
%     \end{choices}
%   \end{question}
% }

% \element{ai}{
%   \begin{question}{Q017}
%     ``Summarize the article in 100 words, focusing on main findings and recommendations." เป็นเทคนิคใด
%     \begin{choices}
%         \correctchoice{Framing}
%         \wrongchoice{Task Specific}
%         \wrongchoice{Bias Mitigation}
%         \wrongchoice{Domain Expertise}
%     \end{choices}
%   \end{question}
% }

% \element{ai}{
%   \begin{question}{Q018}
%     ข้อใดอธิบาย Zero Shot Prompting ได้ถูกต้อง
%     \begin{choices}
%         \wrongchoice{การให้ตัวอย่างหลายตัวก่อนให้โมเดลตอบ}
%         \correctchoice{การให้คำสั่งทั่วไปโดยไม่มีตัวอย่างประกอบ}
%         \wrongchoice{การใช้โมเดลใหม่เสมอทุกครั้ง}
%         \wrongchoice{การฝึกโมเดลด้วยข้อมูลใหม่}
%     \end{choices}
%   \end{question}
% }

% \element{ai}{
%   \begin{question}{Q019}
%     ตัวอย่าง ``Translate ‘Hello’ to Thai" จัดอยู่ในรูปแบบใด
%     \begin{choices}
%         \wrongchoice{Few Shot}
%         \correctchoice{Zero Shot}
%         \wrongchoice{Domain Expertise}
%         \wrongchoice{Contextual Guidance}
%     \end{choices}
%   \end{question}
% }

% \element{ai}{
%   \begin{question}{Q020}
%     ข้อใดเป็นตัวอย่างของ Few Shot Prompting
%     \begin{choices}
%         \wrongchoice{Translate ``Hello" to Thai}
%         \correctchoice{Hello: สวัสดีเจ้า / Good bye: ไปก่อนเน้อเจ้า / Goodnight: …}
%         \wrongchoice{Write a short essay on AI}
%         \wrongchoice{Summarize the article in 100 words}
%     \end{choices}
%   \end{question}
% }

% \element{ai}{
%   \begin{question}{Q021}
%     จุดเด่นของ Few Shot Prompting คืออะไร
%     \begin{choices}
%         \wrongchoice{ทำงานได้โดยไม่ต้องใช้ข้อมูลเลย}
%         \correctchoice{ให้โมเดลเรียนรู้จากตัวอย่างเพียงไม่กี่ตัวอย่าง}
%         \wrongchoice{ใช้ข้อมูลจำนวนมากในการฝึก}
%         \wrongchoice{ทำให้โมเดลไม่ต้องใช้ Context}
%     \end{choices}
%   \end{question}
% }

% \element{ai}{
%   \begin{question}{Q022}
%     หากร้านขายปากกากล่องละ 5 ดอลลาร์ พร้อมโปรซื้อ 1 แถม 1 แล้วซื้อ 10 กล่อง ควรใช้ Pattern ใดในการแก้โจทย์
%     \begin{choices}
%         \correctchoice{ChainofThought Pattern}
%         \wrongchoice{Few Shot Prompting}
%         \wrongchoice{Domain Expertise}
%         \wrongchoice{Bias Mitigation}
%     \end{choices}
%   \end{question}
% }

% \element{ai}{
%   \begin{question}{Q023}
%     ตัวอย่างคำถาม ``How to get a good grade " ใน Tree of Thought ใช้เพื่ออะไร
%     \begin{choices}
%         \wrongchoice{สาธิตการคำนวณเชิงตัวเลข}
%         \correctchoice{สาธิตการวิเคราะห์แบบหลายมุมมอง}
%         \wrongchoice{สาธิตการใช้ ZeroShot Prompting}
%         \wrongchoice{สาธิตการลด Bias}
%     \end{choices}
%   \end{question}
% }

% \element{ai}{
%   \begin{question}{Q024}
%     Management for Sustainability มุ่งเน้นเรื่องใด
%     \begin{choices}
%         \wrongchoice{การจัดเก็บภาษีท้องถิ่น}
%         \correctchoice{การมีแผนและกลไกเพื่อความยั่งยืน และการมีส่วนร่วมของประชาชน}
%         \wrongchoice{การขยายเมืองออกไปสู่ชานเมือง}
%         \wrongchoice{การสร้างสนามกีฬาในเมือง}
%     \end{choices}
%   \end{question}
% }

% \element{ai}{
%   \begin{question}{Q025}
%     Bruce Henderson (1975) เสนอแนวคิดเรื่องใด
%     \begin{choices}
%         \wrongchoice{Value Chain}
%         \correctchoice{Scale curves, Experience curves}
%         \wrongchoice{Five Forces Model}
%         \wrongchoice{Balanced Scorecard}
%     \end{choices}
%   \end{question}
% }

% \element{ai}{
%   \begin{question}{Q026}
%     Michael Porter (1980) เน้นแนวคิดใด
%     \begin{choices}
%         \correctchoice{Value Chain}
%         \wrongchoice{Balanced Scorecard}
%         \wrongchoice{Stakeholder Theory}
%         \wrongchoice{Game Theory}
%     \end{choices}
%   \end{question}
% }

% \element{ai}{
%   \begin{question}{Q027}
%     ข้อใดคือความแตกต่างหลักของแนวคิด Henderson (1975) และ Porter (1980)
%     \begin{choices}
%         \correctchoice{Henderson = เน้นปริมาณ/ประสบการณ์กับต้นทุน, 	Porter = เน้นห่วงโซ่คุณค่าและโครงสร้างธุรกิจ}
%         \wrongchoice{Henderson = เน้น Value Chain, Porter = เน้น Economies of Mass}
%         \wrongchoice{Henderson = เน้น Balanced Scorecard, Porter = เน้น Game Theory}
%         \wrongchoice{ไม่มีความต่าง}
%     \end{choices}
%   \end{question}
% }

\element{ai}{
  \begin{question}{Q028}
    ข้อใด ไม่ใช่ ลักษณะของ Deep Learning (DL)
    \begin{choices}
        \wrongchoice{เป็นสาขาย่อยของ Machine Learning}
        \correctchoice{ใช้โครงข่ายประสาทเทียมหลายชั้น}
        \wrongchoice{ต้องการข้อมูลจำนวนมากในการฝึก}
        \wrongchoice{สามารถทำงานได้โดยไม่ต้องอาศัยข้อมูลใด ๆ}
    \end{choices}
  \end{question}
}

% \element{ai}{
%   \begin{question}{Q029}
%     ผลกระทบต่อแรงงานจาก AI ควรถูกมองอย่างไร
%     \begin{choices}
%         \wrongchoice{เป็นผลบวกเสมอ}
%         \wrongchoice{เป็นผลลบเสมอ}
%         \correctchoice{อาจเป็นได้ทั้งบวกและลบ ขึ้นอยู่กับการจัดการ}
%         \wrongchoice{ไม่มีผลกระทบจริง}
%     \end{choices}
%   \end{question}
% }

% \element{ai}{
%   \begin{question}{Q030}
%     ข้อใดคือแพลตฟอร์มที่เกี่ยวข้องกับ Generative AI ด้านเพลง ในกลุ่ม AI Services
%     \begin{choices}
%         \wrongchoice{ChatGPT.com}
%         \correctchoice{Suno.com}
%         \wrongchoice{Runwayml.com}
%         \wrongchoice{Lovable.dev}
%     \end{choices}
%   \end{question}
% }

% \element{ai}{
%   \begin{question}{Q031}
%     ข้อความ ``a presentation on recent advances in machine learning" จัดอยู่ในส่วนใด
%     \begin{choices}
%         \wrongchoice{Instruction}
%         \correctchoice{Context}
%         \wrongchoice{Input}
%         \wrongchoice{Output}
%     \end{choices}
%   \end{question}
% }

% ======================================================================
% Topic: Big Data
% ======================================================================

\element{bigdata}{
  \begin{question}{Q033}
    Big Data เดิมถูกอธิบายด้วย 3 องค์ประกอบหลัก คืออะไร
    \begin{choices}
        \correctchoice{Volume, Velocity, Variety}
        \wrongchoice{Value, Validity, Veracity}
        \wrongchoice{Visualization, Verification, Variation}
        \wrongchoice{Volume, Value, Veracity}
    \end{choices}
  \end{question}
}

\element{bigdata}{
  \begin{question}{Q034}
    ข้อใดไม่ใช่ความท้าทายของ Big Data ตามนิยาม
    \begin{choices}
        \wrongchoice{Capturing data}
        \wrongchoice{Data storage}
        \wrongchoice{Data analysis}
        \correctchoice{Data cooking}
    \end{choices}
  \end{question}
}

\element{bigdata}{
  \begin{question}{Q035}
    ข้อใดตรงกับ ``Structured Data" ที่สุด
    \begin{choices}
        \wrongchoice{ข้อมูลที่ไม่มีรูปแบบตายตัว}
        \correctchoice{ข้อมูลอยู่ในตาราง รูปแบบ แถว,คอลัมน์ชัดเจน}
        \wrongchoice{ข้อมูลจากไฟล์เสียง}
        \wrongchoice{ข้อมูลรูปภาพ}
    \end{choices}
  \end{question}
}

\element{bigdata}{
  \begin{question}{Q036}
    ข้อใดเป็น Unstructured Data
    \begin{choices}
        \wrongchoice{ตาราง SQL}
        \correctchoice{ข้อความอิสระ, วิดีโอ, ไฟล์เสียง}
        \wrongchoice{ไฟล์ Excel}
        \wrongchoice{ข้อมูลอยู่ในตาราง รูปแบบ แถว,คอลัมน์ชัดเจน}
    \end{choices}
  \end{question}
}

\element{bigdata}{
  \begin{question}{Q037}
    ข้อใด ``ไม่ใช่" ตัวอย่างของ Unstructured Data
    \begin{choices}
        \wrongchoice{รูปภาพ}
        \wrongchoice{ไฟล์ Word}
        \correctchoice{ฐานข้อมูล SQL ตารางพนักงาน}
        \wrongchoice{คลิปเสียง}
    \end{choices}
  \end{question}
}

\element{bigdata}{
  \begin{question}{Q038}
    Metadata คืออะไร
    \begin{choices}
        \wrongchoice{ข้อมูลที่ไม่สำคัญ}
        \correctchoice{ข้อมูลเกี่ยวกับข้อมูล (เช่น วัน, เวลา, พิกัด, กล้องที่ใช้ถ่ายรูป)}
        \wrongchoice{ข้อมูลข้อความ}
        \wrongchoice{ข้อมูลที่อยู่ในรูปแบบ JSON}
    \end{choices}
  \end{question}
}

\element{bigdata}{
  \begin{question}{Q039}
    ความแตกต่างหลักระหว่าง Semistructured กับ Structured คืออะไร
    \begin{choices}
        \correctchoice{Semistructured มีแท็กช่วยกำหนดโครงแต่ไม่เป็นตารางตายตัว ส่วน Structured ต้องเป็นตารางที่ชัดเจน}
        \wrongchoice{Structured เป็นข้อมูลที่เป็นรูปภาพ}
        \wrongchoice{Semistructured ไม่มี metadata}
        \wrongchoice{ไม่มีความต่าง}
    \end{choices}
  \end{question}
}

\element{bigdata}{
  \begin{question}{Q040}
    ข้อใดอธิบาย Volume ได้ถูกต้องที่สุด
    \begin{choices}
        \correctchoice{ปริมาณข้อมูลที่ถูกสร้างและจัดเก็บ}
        \wrongchoice{ความหลากหลายของชนิดข้อมูล}
        \wrongchoice{ความเร็วในการสร้างข้อมูล}
        \wrongchoice{ความถูกต้องแม่นยำของข้อมูล}
    \end{choices}
  \end{question}
}

\element{bigdata}{
  \begin{question}{Q041}
    Velocity ของ Big Data หมายถึง
    \begin{choices}
        \correctchoice{ความเร็วในการสร้างและประมวลผลข้อมูลเพื่อให้ตอบสนองได้แบบ RealTime}
        \wrongchoice{ความถูกต้องของข้อมูล}
        \wrongchoice{ความหลากหลายของไฟล์}
        \wrongchoice{ความแตกต่างระหว่างข้อมูลเชิงคุณภาพกับเชิงปริมาณ}
    \end{choices}
  \end{question}
}

\element{bigdata}{
  \begin{question}{Q042}
    Media vs Social Media ต่างกันอย่างไร
    \begin{choices}
        \correctchoice{Media = สื่อดิจิทัล ส่วน Social Media = เครือข่ายออนไลน์ที่ผู้ใช้สร้าง content}
        \wrongchoice{Media ต้องอยู่ในฐานข้อมูล SQL}
        \wrongchoice{Media เป็นข้อมูลเชิงตัวเลข ส่วน Social Media เป็นข้อมูลตาราง}
        \wrongchoice{ไม่มีความต่าง}
    \end{choices}
  \end{question}
}

\element{bigdata}{
  \begin{question}{Q043}
    ข้อใดคือการเรียงลำดับชั้นของ DIKW ที่ถูกต้อง
    \begin{choices}
        \wrongchoice{Wisdom -> Knowledge -> Information -> Data}
        \correctchoice{Data -> Information -> Knowledge -> Wisdom}
        \wrongchoice{Knowledge -> Wisdom -> Data -> Information}
        \wrongchoice{Information -> Data -> Wisdom -> Knowledge}
    \end{choices}
  \end{question}
}

\element{bigdata}{
  \begin{question}{Q044}
    ถ้าเรามี ``คะแนนสอบนักเรียน" แต่ยังไม่มีการตีความ ถือว่าอยู่ในระดับใดของ DIKW
    \begin{choices}
        \correctchoice{Data}
        \wrongchoice{Information}
        \wrongchoice{Knowledge}
        \wrongchoice{Wisdom}
    \end{choices}
  \end{question}
}

\element{bigdata}{
  \begin{question}{Q045}
    ถ้าเรานำคะแนนสอบไปจัดเรียง ทำตาราง วิเคราะห์เป็นค่าเฉลี่ย จะกลายเป็นระดับใด
    \begin{choices}
        \wrongchoice{Data}
        \correctchoice{Information}
        \wrongchoice{Knowledge}
        \wrongchoice{Wisdom}
    \end{choices}
  \end{question}
}

\element{bigdata}{
  \begin{question}{Q046}
    ถ้าเรานำผลค่าเฉลี่ยไปอธิบายว่านักเรียนกลุ่มหนึ่งมีผลการเรียนอ่อนลงเพราะไม่ได้เตรียมสอบ ถือว่าอยู่ระดับใด
    \begin{choices}
        \wrongchoice{Data}
        \wrongchoice{Information}
        \correctchoice{Knowledge}
        \wrongchoice{Wisdom}
    \end{choices}
  \end{question}
}

\element{bigdata}{
  \begin{question}{Q047}
    ถ้าเรานำความรู้นี้ไปออกนโยบายปรับวิธีสอนเพื่อพัฒนาผลการเรียน ถือว่าอยู่ระดับใด
    \begin{choices}
        \wrongchoice{Data}
        \wrongchoice{Information}
        \wrongchoice{Knowledge}
        \correctchoice{Wisdom}
    \end{choices}
  \end{question}
}

\element{bigdata}{
  \begin{question}{Q048}
    ข้อความ ``Machine learning is a rapidly evolving field that continues to transform industries by enabling systems to learn from data" จัดอยู่ในส่วนใด
    \begin{choices}
        \wrongchoice{Instruction}
        \wrongchoice{Context}
        \correctchoice{Input}
        \wrongchoice{Output}
    \end{choices}
  \end{question}
}

\element{bigdata}{
  \begin{question}{Q050}
    Machine Learning (ML) แตกต่างจาก AI ทั่วไปอย่างไร
    \begin{choices}
        \wrongchoice{Machine Learning คือการเขียนโปรแกรมให้ทำงานแบบกำหนดกฎตายตัว}
        \wrongchoice{Machine Learning คืออัลกอริทึมที่ปรับปรุงประสิทธิภาพได้เองเมื่อเจอข้อมูลมากขึ้น}
        \wrongchoice{Machine Learning เป็นระบบที่คิดเชิงเหตุผลเหมือนมนุษย์}
        \correctchoice{Machine Learning เน้นใช้ Big Data อย่างเดียว}
    \end{choices}
  \end{question}
}

\element{bigdata}{
  \begin{question}{Q051}
    ชั้น Business Layer เกี่ยวข้องกับเรื่องใด
    \begin{choices}
        \correctchoice{โมเดลธุรกิจ ความเป็นส่วนตัว และการจัดการระบบโดยรวม}
        \wrongchoice{เฉพาะการเชื่อมต่อ Bluetooth}
        \wrongchoice{การตรวจจับข้อมูลสภาพอากาศ}
        \wrongchoice{การประมวลผล Big Data โดยตรง}
    \end{choices}
  \end{question}
}

\element{bigdata}{
  \begin{question}{Q052}
    ไฟล์ Excel, CSV, Word, JSON จะถูกจัดในหมวดใด
    \begin{choices}
        \wrongchoice{Media}
        \correctchoice{Docs}
        \wrongchoice{Data Storage}
        \wrongchoice{Social Media}
    \end{choices}
  \end{question}
}

% ======================================================================
% Topic: Cloud Computing
% ======================================================================

\element{cloudcomputing}{
  \begin{question}{Q053}
    Variety ใน Big Data หมายถึงอะไร
    \begin{choices}
        \wrongchoice{ความเร็วของข้อมูล}
        \correctchoice{ชนิดและธรรมชาติของข้อมูลที่หลากหลาย เช่น ข้อความ รูปภาพ วิดีโอ เสียง}
        \wrongchoice{ปริมาณการจัดเก็บ}
        \wrongchoice{การใช้ cloud computing}
    \end{choices}
  \end{question}
}

\element{cloudcomputing}{
  \begin{question}{Q054}
    Cloud ให้บริการผ่านเครือข่ายแบบใดได้บ้าง
    \begin{choices}
        \wrongchoice{ได้เฉพาะอินเทอร์เน็ตบ้านเท่านั้น}
        \correctchoice{ทั้งเครือข่ายสาธารณะและเครือข่ายส่วนตัว เช่น WAN, LAN, VPN}
        \wrongchoice{ได้เฉพาะบลูทูธ}
        \wrongchoice{ต้องต่อสายตรงระหว่างเครื่องเท่านั้น}
    \end{choices}
  \end{question}
}

\element{cloudcomputing}{
  \begin{question}{Q055}
    ข้อใดคือความหมายของ Public Cloud
    \begin{choices}
        \wrongchoice{ใช้เฉพาะในองค์กรเดียว}
        \correctchoice{มีพื้นที่มากและขยายได้ง่าย เหมาะกับงานพัฒนาและทำงานร่วมกัน}
        \wrongchoice{ใช้หลายองค์กรร่วมกันเฉพาะในอุตสาหกรรมเดียว}
        \wrongchoice{ผสมผสานระหว่างระบบออฟไลน์กับออนไลน์}
    \end{choices}
  \end{question}
}

\element{cloudcomputing}{
  \begin{question}{Q056}
    Private Cloud เหมาะสมกับสถานการณ์ใดมากที่สุด
    \begin{choices}
        \wrongchoice{เกมออนไลน์ขนาดใหญ่}
        \correctchoice{ธุรกิจที่ต้องการความปลอดภัยสูง}
        \wrongchoice{การแชร์ข้อมูลระหว่างหลายองค์กรในอุตสาหกรรมเดียวกัน}
        \wrongchoice{นักเรียนที่อยากใช้พื้นที่ฟรี}
    \end{choices}
  \end{question}
}

\element{cloudcomputing}{
  \begin{question}{Q057}
    Community Cloud หมายถึงข้อใด
    \begin{choices}
        \wrongchoice{ใช้แค่ภายในบ้าน}
        \correctchoice{แพลตฟอร์มที่หลายองค์กรในอุตสาหกรรมเดียวกันใช้ร่วมกัน}
        \wrongchoice{การเก็บข้อมูลบนมือถือ}
        \wrongchoice{ใช้พื้นที่จาก USB Drive}
    \end{choices}
  \end{question}
}

\element{cloudcomputing}{
  \begin{question}{Q058}
    (Infrastructure as a Service) IaaS คือบริการแบบใด
    \begin{choices}
        \correctchoice{ให้ผู้ใช้เช่าใช้ฮาร์ดแวร์และระบบพื้นฐานแบบตามต้องการ}
        \wrongchoice{ให้ผู้ใช้เฉพาะซอฟต์แวร์สำเร็จรูป}
        \wrongchoice{ให้ผู้ใช้ซื้อเครื่องเซิร์ฟเวอร์เอง}
        \wrongchoice{ให้ผู้ใช้เชื่อมต่อ WiFi ฟรี}
    \end{choices}
  \end{question}
}

\element{cloudcomputing}{
  \begin{question}{Q059}
    (Platform as a Service) PaaS ให้บริการอะไรเป็นหลัก
    \begin{choices}
        \wrongchoice{ซอฟต์แวร์สำเร็จรูปพร้อมใช้งานทันที}
        \correctchoice{เครื่องมือสำหรับพัฒนาและรันแอปพลิเคชัน}
        \wrongchoice{เครื่องคอมพิวเตอร์จริงที่ต้องซื้อเอง}
        \wrongchoice{การเชื่อมต่ออินเทอร์เน็ตฟรี}
    \end{choices}
  \end{question}
}

\element{cloudcomputing}{
  \begin{question}{Q060}
    (Software as a Service) SaaS คืออะไร
    \begin{choices}
        \wrongchoice{การเช่าเครื่องคอมพิวเตอร์จริงจากผู้ให้บริการ}
        \correctchoice{การใช้ซอฟต์แวร์เป็นบริการผ่านอินเทอร์เน็ต}
        \wrongchoice{การซื้อโปรแกรมแล้วติดตั้งลงเครื่อง}
        \wrongchoice{การพัฒนาแพลตฟอร์มเองทั้งหมด}
    \end{choices}
  \end{question}
}

\element{cloudcomputing}{
  \begin{question}{Q061}
    บริการแบบ (Infrastructure as a Service) IaaS มักถูกใช้งานโดยใคร
    \begin{choices}
        \wrongchoice{ผู้ใช้งานทั่วไป (End users)}
        \wrongchoice{นักพัฒนาซอฟต์แวร์ (Software developers)}
        \correctchoice{ผู้ดูแลระบบ IT}
        \wrongchoice{นักเรียนมัธยม}
    \end{choices}
  \end{question}
}

\element{cloudcomputing}{
  \begin{question}{Q062}
    จุดเด่นสำคัญของ Edge Computing คืออะไร
    \begin{choices}
        \correctchoice{ช่วยลดเวลาในการตอบสนอง (response time)}
        \wrongchoice{ทำให้การประมวลผลช้าลง}
        \wrongchoice{ใช้ได้เฉพาะในเกมเท่านั้น}
        \wrongchoice{ต้องมีอินเทอร์เน็ตตลอดเวลา}
    \end{choices}
  \end{question}
}

\element{cloudcomputing}{
  \begin{question}{Q063}
    ข้อใดต่อไปนี้ไม่ใช่ตัวอย่างของ AI as a service (AIaaS)
    \begin{choices}
        \wrongchoice{OpenAI}
        \wrongchoice{Google AI}
        \correctchoice{Canva}
        \wrongchoice{Amazon Web Services}
    \end{choices}
  \end{question}
}

\element{cloudcomputing}{
  \begin{question}{Q064}
    คำว่า ``Cloud technology" หมายถึงอะไร
    \begin{choices}
        \wrongchoice{ฮาร์ดดิสก์ในเครื่องคอมพิวเตอร์}
        \correctchoice{เครือข่ายหรืออินเทอร์เน็ต}
        \wrongchoice{โปรแกรมที่ติดตั้งในเครื่อง}
        \wrongchoice{แฟลชไดรฟ์ที่เสียบคอมพิวเตอร์}
    \end{choices}
  \end{question}
}

\element{cloudcomputing}{
  \begin{question}{Q065}
    Hybrid Cloud มีจุดเด่นคืออะไร
    \begin{choices}
        \correctchoice{ผสมการใช้ Public และ Private Cloud เข้าด้วยกัน}
        \wrongchoice{ใช้เฉพาะกับการประชุมออนไลน์}
        \wrongchoice{เก็บข้อมูลได้เฉพาะในฮาร์ดดิสก์ภายในเครื่อง}
        \wrongchoice{ใช้สำหรับส่งอีเมลเท่านั้น}
    \end{choices}
  \end{question}
}

\element{cloudcomputing}{
  \begin{question}{Q066}
    ตัวอย่างบริการใดตรงกับ SaaS
    \begin{choices}
        \correctchoice{Gmail, Trello, Office 365}
        \wrongchoice{AWS, Stackscale}
        \wrongchoice{Heroku, Cloud Foundry}
        \wrongchoice{VMware}
    \end{choices}
  \end{question}
}

% ======================================================================
% Topic: General
% ======================================================================

\element{general}{
  \begin{question}{Q067}
    Economies of Scale มีผลต่อธุรกิจอย่างไร
    \begin{choices}
        \wrongchoice{ทำให้ต้นทุนต่อหน่วยสูงขึ้นเมื่อผลิตเพิ่ม}
        \correctchoice{ทำให้ต้นทุนต่อหน่วยลดลงเมื่อผลิตเพิ่ม}
        \wrongchoice{ไม่มีผลต่อการผลิต}
        \wrongchoice{ทำให้ต้นทุนเพิ่มคงที่}
    \end{choices}
  \end{question}
}

\element{general}{
  \begin{question}{Q068}
    ข้อใดคือ Internal Economies of Scale
    \begin{choices}
        \wrongchoice{ได้รับสิทธิประโยชน์ทางภาษีจากรัฐบาล}
        \correctchoice{การซื้อวัตถุดิบปริมาณมากทำให้ราคาถูกลง}
        \wrongchoice{การสนับสนุนจากสมาคมอุตสาหกรรม}
        \wrongchoice{การช่วยเหลือจากต่างประเทศ}
    \end{choices}
  \end{question}
}

\element{general}{
  \begin{question}{Q069}
    ข้อใดคือ External Economies of Scale
    \begin{choices}
        \wrongchoice{การซื้อวัตถุดิบจำนวนมากโดยบริษัท}
        \correctchoice{การได้รับ preferential treatment จากรัฐบาล}
        \wrongchoice{การจัดการสายการผลิตภายในบริษัท}
        \wrongchoice{การรวมแผนกเพื่อประหยัดค่าใช้จ่าย}
    \end{choices}
  \end{question}
}

\element{general}{
  \begin{question}{Q070}
    Diseconomies of Scale หมายถึงอะไร
    \begin{choices}
        \wrongchoice{การที่ขนาดเล็กเกินไปทำให้ไม่คุ้มทุน}
        \correctchoice{การที่บริษัทขยายใหญ่เกินไปจนประสิทธิภาพลดลงและต้นทุนสูงขึ้น}
        \wrongchoice{การแข่งขันทำให้ราคาตกต่ำ}
        \wrongchoice{การที่บริษัทไม่ได้รับการสนับสนุนจากรัฐบาล}
    \end{choices}
  \end{question}
}

\element{general}{
  \begin{question}{Q071}
    ตาม กฎของมัวร์ (Moore’s Law) กล่าวไว้อย่างไร
    \begin{choices}
        \wrongchoice{จำนวนทรานซิสเตอร์จะลดลงครึ่งหนึ่งทุก 2 ปี}
        \correctchoice{จำนวนทรานซิสเตอร์บนไมโครชิปจะเพิ่มขึ้น 2 เท่าทุก 2 ปี}
        \wrongchoice{ความเร็วของคอมพิวเตอร์จะลดลงทุก 5 ปี}
        \wrongchoice{ราคาคอมพิวเตอร์จะคงที่ตลอดเวลา}
    \end{choices}
  \end{question}
}

\element{general}{
  \begin{question}{Q072}
    ``ทําตัวเปนเชฟอาหารระดับโลก เดียวฉันจะบรรยายอาหารทีกิน แล้วคุณคอมเมนต์เกียวกับเมนู นันให้หน่อย" เป็นตัวอย่างของอะไร
    \begin{choices}
        \correctchoice{Persona Pattern}
        \wrongchoice{Audience Persona Pattern}
        \wrongchoice{Template Pattern}
        \wrongchoice{Interview Pattern}
    \end{choices}
  \end{question}
}

\element{general}{
  \begin{question}{Q073}
    อธิบายการทํางานของห่วงโซ่อุปทานในร้านขายของชําของสหรัฐฯ ให้ฉันฟง สมมติว่าฉันคือ เจงกิสข่าน เป็นตัวอย่างของอะไร
    \begin{choices}
        \wrongchoice{Persona Pattern}
        \correctchoice{Audience Persona Pattern}
        \wrongchoice{Template Pattern}
        \wrongchoice{Interview Pattern}
    \end{choices}
  \end{question}
}

\element{general}{
  \begin{question}{Q074}
    สร้างโปรแกรมออกกําลังกายตามเทมเพลตนี: ชือท่า, จํานวนครัง x เซ็ต, กลุ่มกล้ามเนือทีใช้, ระดับ ความยาก 15, หมายเหตุเรืองท่าทาง เป็นตัวอย่างของอะไร
    \begin{choices}
        \wrongchoice{Persona Pattern}
        \wrongchoice{Audience Persona Pattern}
        \correctchoice{Template Pattern}
        \wrongchoice{Interview Pattern}
    \end{choices}
  \end{question}
}

\element{general}{
  \begin{question}{Q075}
    ถ้าเป็นเครือข่ายเล็ก ๆ ในสำนักงานทั่วไปที่มีอุปกรณ์ไม่กี่ร้อยเครื่อง ควรใช้แบบใด
    \begin{choices}
        \correctchoice{IPv4}
        \wrongchoice{IPv6}
        \wrongchoice{ทั้งคู่ใช้ไม่ได้}
        \wrongchoice{ต้องใช้เครือข่ายพิเศษเท่านั้น}
    \end{choices}
  \end{question}
}

\element{general}{
  \begin{question}{Q076}
    โครงสร้างพื้นฐานแบบ Semiintelligent มีลักษณะอย่างไร
    \begin{choices}
        \wrongchoice{เก็บข้อมูลและตัดสินใจเองได้}
        \correctchoice{เก็บข้อมูลแต่ไม่สามารถตัดสินใจจากข้อมูลได้}
        \wrongchoice{ไม่เก็บข้อมูลใด ๆ}
        \wrongchoice{ตัดสินใจอัตโนมัติโดยไม่ต้องใช้มนุษย์}
    \end{choices}
  \end{question}
}

\element{general}{
  \begin{question}{Q077}
    แนวคิด Scale curve และ Experience curve สื่อถึงอะไร
    \begin{choices}
        \correctchoice{การผลิตเพิ่มขึ้นทำให้ต้นทุนต่อหน่วยลดลง}
        \wrongchoice{การผลิตมากขึ้นทำให้ต้นทุนสูงขึ้น}
        \wrongchoice{ต้นทุนไม่เปลี่ยนแปลงตามปริมาณ}
        \wrongchoice{การลงทุนไม่เกี่ยวกับต้นทุน}
    \end{choices}
  \end{question}
}

\element{general}{
  \begin{question}{Q078}
    Economies of Scale หมายถึงอะไร
    \begin{choices}
        \wrongchoice{ต้นทุนต่อหน่วยเพิ่มขึ้นเมื่อผลิตมากขึ้น}
        \correctchoice{ข้อได้เปรียบที่บริษัทขนาดใหญ่มีจากการลดต้นทุนเมื่อปริมาณเพิ่ม}
        \wrongchoice{การได้รับเงินอุดหนุนจากรัฐบาลเท่านั้น}
        \wrongchoice{การผลิตน้อยแต่คุณภาพสูง}
    \end{choices}
  \end{question}
}

\element{general}{
  \begin{question}{Q079}
    Psychographic segmentation เน้นปัจจัยใด
    \begin{choices}
        \wrongchoice{อายุ เพศ รายได้}
        \correctchoice{บุคลิกภาพ ไลฟ์สไตล์ ความเชื่อ}
        \wrongchoice{ที่อยู่ ภูมิศาสตร์}
        \wrongchoice{ความถี่การซื้อสินค้า}
    \end{choices}
  \end{question}
}

\element{general}{
  \begin{question}{Q080}
    ชั้น Perception Layer มีหน้าที่หลักคืออะไร
    \begin{choices}
        \wrongchoice{วิเคราะห์ข้อมูลจำนวนมาก}
        \correctchoice{ตรวจจับและเก็บข้อมูลจากสิ่งแวดล้อมด้วยเซนเซอร์}
        \wrongchoice{จัดการเรื่องธุรกิจและความเป็นส่วนตัว}
        \wrongchoice{ส่งต่อข้อมูลผ่านเครือข่าย}
    \end{choices}
  \end{question}
}

\element{general}{
  \begin{question}{Q081}
    ชั้น Transport Layer ทำหน้าที่ใด
    \begin{choices}
        \correctchoice{ส่งข้อมูลจากเซนเซอร์ไปยังชั้นอื่น ๆ ผ่านเครือข่าย}
        \wrongchoice{ให้บริการเฉพาะเจาะจงกับผู้ใช้}
        \wrongchoice{จัดการข้อมูลจำนวนมากบนคลาวด์}
        \wrongchoice{ควบคุมโมเดลธุรกิจและการใช้งานระบบ}
    \end{choices}
  \end{question}
}

\element{general}{
  \begin{question}{Q082}
    ชั้น Application Layer มีหน้าที่ใด
    \begin{choices}
        \wrongchoice{ตรวจจับสภาพแวดล้อม}
        \correctchoice{ให้บริการแอปพลิเคชันเฉพาะสำหรับผู้ใช้}
        \wrongchoice{ควบคุมเครือข่ายการเชื่อมต่อ}
        \wrongchoice{เก็บข้อมูลระยะยาว}
    \end{choices}
  \end{question}
}

\element{general}{
  \begin{question}{Q083}
    การเพิ่มนวัตกรรม (
    \begin{choices}
        \correctchoice{ทำให้เส้น Output (Y) ขยับขึ้น}
        \wrongchoice{ทำให้ทุน (K) ลดลง}
        \wrongchoice{ทำให้แรงงาน (eL) หายไป}
        \wrongchoice{ทำให้ค่าเสื่อมราคามากขึ้น}
    \end{choices}
  \end{question}
}

\element{general}{
  \begin{question}{Q084}
    Automation ช่วยเมืองอัจฉริยะในด้านใด
    \begin{choices}
        \correctchoice{ทำให้ระบบสามารถวิเคราะห์ข้อมูลจากเซนเซอร์และสั่งการได้เอง}
        \wrongchoice{ทำให้เมืองมีระบบไฟถนนเพิ่มขึ้น}
        \wrongchoice{ทำให้เมืองสร้างอาคารสูงขึ้น}
        \wrongchoice{ทำให้ทุกคนต้องควบคุมเองทั้งหมด}
    \end{choices}
  \end{question}
}

\element{general}{
  \begin{question}{Q085}
    ข้อใดเป็นตัวอย่าง application ที่ ``มักรันบนคลาวด์"
    \begin{choices}
        \wrongchoice{โปรแกรมเครื่องคิดเลขออฟไลน์}
        \correctchoice{การส่งอีเมลและประชุมออนไลน์}
        \wrongchoice{เกมออฟไลน์บนเครื่องเดียว}
        \wrongchoice{โปรแกรมจดโน้ตแบบไม่เชื่อมเน็ต}
    \end{choices}
  \end{question}
}

% ======================================================================
% Topic: Iot
% ======================================================================

\element{iot}{
  \begin{question}{Q086}
    IoT Gateway ทำหน้าที่อะไร
    \begin{choices}
        \wrongchoice{เป็นที่เก็บข้อมูลระยะยาว}
        \correctchoice{เป็นตัวเชื่อมต่ออุปกรณ์ IoT เข้าสู่เครือข่ายและคลาวด์}
        \wrongchoice{เป็นหน้าจอให้ผู้ใช้ควบคุม}
        \wrongchoice{เป็นอุปกรณ์ตรวจจับอุณหภูมิ}
    \end{choices}
  \end{question}
}

\element{iot}{
  \begin{question}{Q087}
    ``Internet of Things with IP address" หมายถึงอะไร
    \begin{choices}
        \correctchoice{อุปกรณ์ที่เชื่อมต่ออินเทอร์เน็ตได้เองโดยไม่ต้องพึ่งคน}
        \wrongchoice{อุปกรณ์ที่ต้องมีคนกดปุ่มทุกครั้ง}
        \wrongchoice{ระบบที่ใช้ได้เฉพาะคอมพิวเตอร์}
        \wrongchoice{อุปกรณ์ที่ไม่มีเซนเซอร์}
    \end{choices}
  \end{question}
}

\element{iot}{
  \begin{question}{Q088}
    Other ID and Sensor systems (IoT without IP address) หมายถึงอะไร
    \begin{choices}
        \wrongchoice{ระบบที่ไม่ใช้เซนเซอร์เลย}
        \correctchoice{ระบบที่มีเซนเซอร์ แต่ไม่เชื่อมต่อโดยตรงกับอินเทอร์เน็ต}
        \wrongchoice{ระบบที่มีเฉพาะมือถือและแท็บเล็ต}
        \wrongchoice{เครือข่ายที่ใช้เฉพาะ WiFi เท่านั้น}
    \end{choices}
  \end{question}
}

\element{iot}{
  \begin{question}{Q089}
    ถ้าต้องรองรับอุปกรณ์ IoT หลายพันล้านชิ้นทั่วโลก ควรใช้แบบใด
    \begin{choices}
        \wrongchoice{IPv4}
        \correctchoice{IPv6}
        \wrongchoice{ทั้งคู่ใช้ไม่ได้}
        \wrongchoice{ต้องใช้เครือข่ายพิเศษเท่านั้น}
    \end{choices}
  \end{question}
}

\element{iot}{
  \begin{question}{Q090}
    เซนเซอร์ใดใช้ตรวจสอบทิศทางและการหมุน
    \begin{choices}
        \correctchoice{Gyroscope Sensor}
        \wrongchoice{Proximity Sensor}
        \wrongchoice{Rain Sensor}
        \wrongchoice{Alcohol Sensor}
    \end{choices}
  \end{question}
}

\element{iot}{
  \begin{question}{Q091}
    Digital Infrastructure ใน Smart Infrastructure มีอะไรบ้าง
    \begin{choices}
        \wrongchoice{รถไฟฟ้า, ถนน, เขื่อน}
        \correctchoice{Sensors, IoT, Networks, Big Data, Machine Learning}
        \wrongchoice{โรงเรียน, โรงพยาบาล, ตลาด}
        \wrongchoice{แม่น้ำและภูเขา}
    \end{choices}
  \end{question}
}

\element{iot}{
  \begin{question}{Q092}
    Cloud ในระบบ IoT มีไว้ทำอะไร
    \begin{choices}
        \wrongchoice{แจกสัญญาณ WiFi}
        \correctchoice{วิเคราะห์และจัดเก็บข้อมูลจากอุปกรณ์}
        \wrongchoice{ตรวจจับควันโดยตรง}
        \wrongchoice{เป็นที่เก็บข้อมูลระยะยาว}
    \end{choices}
  \end{question}
}

\element{iot}{
  \begin{question}{Q093}
    Processing Layer ทำอะไร
    \begin{choices}
        \wrongchoice{เก็บข้อมูลบนอุปกรณ์ IoT โดยตรง}
        \correctchoice{วิเคราะห์ ประมวลผล และจัดเก็บข้อมูลจำนวนมหาศาล}
        \wrongchoice{ส่งข้อมูลต่อด้วยสัญญาณไร้สาย}
        \wrongchoice{แสดงผลลัพธ์ให้ผู้ใช้}
    \end{choices}
  \end{question}
}

\element{iot}{
  \begin{question}{Q094}
    Sensors และ IoT (Internet of Things) ช่วยเมืองอัจฉริยะอย่างไร
    \begin{choices}
        \correctchoice{เก็บข้อมูลสิ่งแวดล้อม การจราจร คุณภาพอากาศ และเชื่อมต่ออุปกรณ์ต่าง ๆ}
        \wrongchoice{ทำให้ประชาชนใช้โซเชียลมีเดียได้มากขึ้น}
        \wrongchoice{สร้างความบันเทิงเท่านั้น}
        \wrongchoice{ใช้แทนการวางผังเมืองทั้งหมด}
    \end{choices}
  \end{question}
}

\element{iot}{
  \begin{question}{Q095}
    ข้อใดเป็นตัวอย่าง Sensor Data
    \begin{choices}
        \wrongchoice{Metadata ของรูปถ่าย}
        \correctchoice{ข้อมูลจาก Smart electric meters}
        \wrongchoice{Clickstream logs}
        \wrongchoice{Document repository}
    \end{choices}
  \end{question}
}

\element{iot}{
  \begin{question}{Q096}
    เหตุผลที่ทำให้ IoT มีความสามารถในการประมวลผลและวิเคราะห์มากขึ้นคืออะไร
    \begin{choices}
        \correctchoice{การมีการเชื่อมต่อเซนเซอร์และอุปกรณ์เข้ากับอินเทอร์เน็ต}
        \wrongchoice{การตัดการเชื่อมต่อจากระบบเครือข่าย}
        \wrongchoice{การใช้พลังงานจากมนุษย์แทนเครื่องจักร}
        \wrongchoice{การลดจำนวนอุปกรณ์ที่เชื่อมต่อ}
    \end{choices}
  \end{question}
}

\element{iot}{
  \begin{question}{Q097}
    ถ้าต้องการตรวจจับระดับแอลกอฮอล์ในลมหายใจ ควรใช้เซนเซอร์ใด
    \begin{choices}
        \wrongchoice{Gas Sensor}
        \correctchoice{Alcohol Sensor}
        \wrongchoice{PIR Sensor}
        \wrongchoice{Ultrasonic Sensor}
    \end{choices}
  \end{question}
}

% ======================================================================
% Topic: Metaverse
% ======================================================================

\element{metaverse}{
  \begin{question}{Q098}
    คำว่า BlueCollar หมายถึงอะไร
    \begin{choices}
        \wrongchoice{งานที่เกี่ยวข้องกับการทำงานในสำนักงาน}
        \wrongchoice{งานด้านวิชาชีพ เช่น ทนาย แพทย์}
        \correctchoice{งานใช้แรงงานหรือทักษะ เช่น ช่างไฟฟ้า ช่างยนต์}
        \wrongchoice{งานที่เกี่ยวข้องกับการบริหารจัดการ}
    \end{choices}
  \end{question}
}

\element{metaverse}{
  \begin{question}{Q099}
    คำว่า WhiteCollar หมายถึงอะไร
    \begin{choices}
        \wrongchoice{งานภาคการผลิต}
        \correctchoice{งานที่ทำในออฟฟิศ เช่น ทนาย แพทย์ หรือพนักงานธนาคาร}
        \wrongchoice{งานใช้แรงงานก่อสร้าง}
        \wrongchoice{งานที่ไม่ใช้คอมพิวเตอร์}
    \end{choices}
  \end{question}
}

\element{metaverse}{
  \begin{question}{Q100}
    ข้อความ ``a summary of no more than 150 words, highlighting the key advances mentioned." จัดอยู่ในส่วนใด
    \begin{choices}
        \wrongchoice{Instruction}
        \wrongchoice{Context}
        \wrongchoice{Input}
        \correctchoice{Output}
    \end{choices}
  \end{question}
}

\element{metaverse}{
  \begin{question}{Q101}
    จุดประสงค์หลักของ RetrievalAugmented Generation (RAG) คืออะไร
    \begin{choices}
        \wrongchoice{ทำให้โมเดลสร้างเนื้อหาจากศูนย์}
        \correctchoice{ทำให้โมเดลเข้าถึงเอกสารที่กำหนดไว้อย่างแม่นยำ}
        \wrongchoice{ทำให้โมเดลลดขนาด parameter ลง}
        \wrongchoice{ทำให้โมเดลเรียนรู้แบบ unsupervised}
    \end{choices}
  \end{question}
}

\element{metaverse}{
  \begin{question}{Q102}
    ``Write a short paragraph on New York City, highlighting its iconic landmarks." จัดเป็นเทคนิคใด
    \begin{choices}
        \correctchoice{Contextual Guidance}
        \wrongchoice{Task Specific}
        \wrongchoice{User Feedback Loop}
        \wrongchoice{Bias Mitigation}
    \end{choices}
  \end{question}
}

\element{metaverse}{
  \begin{question}{Q103}
    ตัวเลือกใด ``จับคู่ผิด"
    \begin{choices}
        \wrongchoice{Library search = Old Fashion}
        \wrongchoice{Trade publication extraction = Old Fashion}
        \correctchoice{Geolocation stream = Old Fashion}
        \wrongchoice{Social networks API = Today}
    \end{choices}
  \end{question}
}

\element{metaverse}{
  \begin{question}{Q104}
    จุดประสงค์หลักของ RAG คืออะไร
    \begin{choices}
        \wrongchoice{ทำให้โมเดลสร้างเนื้อหาจากศูนย์}
        \correctchoice{ทำให้โมเดลเข้าถึงเอกสารที่กำหนดไว้อย่างแม่นยำ}
        \wrongchoice{ทำให้โมเดลลดขนาด parameter ลง}
        \wrongchoice{ทำให้โมเดลเรียนรู้แบบ unsupervised}
    \end{choices}
  \end{question}
}

% ======================================================================
% Topic: Smart City
% ======================================================================

\element{smartcity}{
  \begin{question}{Q105}
    Smart Citizens ในเมืองอัจฉริยะหมายถึงอะไร
    \begin{choices}
        \correctchoice{พลเมืองที่มีความรู้ เท่าทันเทคโนโลยี และมีส่วนร่วมในการพัฒนาเมือง}
        \wrongchoice{ประชาชนที่ใช้สมาร์ทโฟนทุกคน}
        \wrongchoice{ประชาชนที่อยู่ในเมืองใหญ่เท่านั้น}
        \wrongchoice{ประชาชนที่ทำงานด้านไอทีเท่านั้น}
    \end{choices}
  \end{question}
}

\element{smartcity}{
  \begin{question}{Q106}
    Physical Infrastructure ที่เป็นส่วนหนึ่งของ Smart Infrastructure ได้แก่ข้อใด
    \begin{choices}
        \correctchoice{พลังงาน น้ำ การขนส่ง ขยะ และโทรคมนาคม}
        \wrongchoice{เซนเซอร์ ปัญญาประดิษฐ์ คลาวด์ และ Big Data}
        \wrongchoice{สื่อสังคมออนไลน์ และแอปมือถือ}
        \wrongchoice{บล็อกเชน และโทเคน}
    \end{choices}
  \end{question}
}

\element{smartcity}{
  \begin{question}{Q107}
    Intelligent infrastructure ทำหน้าที่อย่างไร
    \begin{choices}
        \correctchoice{เก็บข้อมูลและนำเสนอเพื่อช่วยมนุษย์ตัดสินใจ}
        \wrongchoice{ไม่เก็บข้อมูลใด ๆ}
        \wrongchoice{ตัดสินใจได้เองโดยไม่ต้องมีมนุษย์}
        \wrongchoice{ใช้ข้อมูลเพียงเพื่อการเก็บสถิติเท่านั้น}
    \end{choices}
  \end{question}
}

\element{smartcity}{
  \begin{question}{Q108}
    Smart infrastructure แตกต่างจากประเภทอื่นตรงไหน
    \begin{choices}
        \wrongchoice{เก็บข้อมูลแต่ไม่ตัดสินใจ}
        \wrongchoice{ต้องใช้คนตัดสินใจเสมอ}
        \correctchoice{ตัดสินใจและปรับตัวได้เองแบบอัตโนมัติ โดยไม่ต้องมีมนุษย์}
        \wrongchoice{ใช้เทคโนโลยีแต่ไม่สามารถเปลี่ยนแปลงตามสภาพจริง}
    \end{choices}
  \end{question}
}

\element{smartcity}{
  \begin{question}{Q109}
    Smart Environment มีเป้าหมายหลักคืออะไร
    \begin{choices}
        \wrongchoice{การสร้างตึกสูงในเมือง}
        \correctchoice{การรักษาสมดุลธรรมชาติและใช้ทรัพยากรอย่างยั่งยืน}
        \wrongchoice{การสร้างสนามบินใหม่}
        \wrongchoice{การเพิ่มปริมาณยานพาหนะในเมือง}
    \end{choices}
  \end{question}
}

\element{smartcity}{
  \begin{question}{Q110}
    ข้อใดคือ ตัวอย่างของ Smart Energy
    \begin{choices}
        \wrongchoice{การสร้างสวนสาธารณะ}
        \wrongchoice{ระบบขนส่งสาธารณะ}
        \correctchoice{พลังงานหมุนเวียน และระบบ Smart Grid}
        \wrongchoice{การออกแบบเมืองเพื่อการท่องเที่ยว}
    \end{choices}
  \end{question}
}

\element{smartcity}{
  \begin{question}{Q111}
    Smart Mobility เน้นในเรื่องใดมากที่สุด
    \begin{choices}
        \correctchoice{การเดินทางและระบบขนส่งที่สะดวก ปลอดภัย และเป็นมิตรต่อสิ่งแวดล้อม}
        \wrongchoice{การสร้างตลาดและแหล่งชุมชนใหม่}
        \wrongchoice{การใช้พลังงานไฟฟ้าในครัวเรือน}
        \wrongchoice{การสร้างสวนดอกไม้}
    \end{choices}
  \end{question}
}

\element{smartcity}{
  \begin{question}{Q112}
    ข้อใด ไม่ใช่ หลักการของ Smart Living
    \begin{choices}
        \wrongchoice{Smart Health (สุขภาพและสุขอนามัยของคน)}
        \wrongchoice{Public Safety & Security (ความปลอดภัยจากอาชญากรรมและภัยพิบัติ)}
        \wrongchoice{Smart Built Environment (สิ่งแวดล้อมการอยู่อาศัยที่ชาญฉลาด)}
        \correctchoice{Smart Agriculture (เกษตรอัจฉริยะ)}
    \end{choices}
  \end{question}
}

\element{smartcity}{
  \begin{question}{Q113}
    Smart Environment มุ่งเน้นด้านใดมากที่สุด
    \begin{choices}
        \wrongchoice{การพัฒนาเทคโนโลยีสารสนเทศ}
        \correctchoice{การรักษาสมดุลธรรมชาติและการใช้ทรัพยากรอย่างยั่งยืน}
        \wrongchoice{การสร้างระบบไฟฟ้าพลังงานใหม่}
        \wrongchoice{การพัฒนาเขตเศรษฐกิจพิเศษ}
    \end{choices}
  \end{question}
}

\element{smartcity}{
  \begin{question}{Q114}
    Smart People มุ่งพัฒนาเรื่องใด
    \begin{choices}
        \correctchoice{การใช้เทคโนโลยีดิจิทัลเพื่อเรียนรู้ตลอดชีวิตและลดความเหลื่อมล้ำ}
        \wrongchoice{การสร้างพลังงานไฟฟ้าใหม่}
        \wrongchoice{การจัดการระบบขนส่ง}
        \wrongchoice{การพัฒนาภาคเกษตรกรรม}
    \end{choices}
  \end{question}
}

\element{smartcity}{
  \begin{question}{Q115}
    Smart Governance เกี่ยวข้องกับอะไร
    \begin{choices}
        \correctchoice{การบริหารภาครัฐที่โปร่งใส และการเข้าถึงบริการภาครัฐ}
        \wrongchoice{การสร้างพลังงานหมุนเวียน}
        \wrongchoice{การพัฒนาเมืองท่องเที่ยว}
        \wrongchoice{การสร้างระบบไฟฟ้าแรงสูง}
    \end{choices}
  \end{question}
}

\element{smartcity}{
  \begin{question}{Q116}
    Smart Economy มีเป้าหมายเพื่ออะไร
    \begin{choices}
        \wrongchoice{การสร้างสวนสาธารณะในเมือง}
        \correctchoice{การส่งเสริมธุรกิจใหม่ ๆ และระบบเศรษฐกิจดิจิทัล}
        \wrongchoice{การพัฒนาพลังงานไฟฟ้า}
        \wrongchoice{การพัฒนาภาคเกษตรกรรมแบบดั้งเดิม}
    \end{choices}
  \end{question}
}

\element{smartcity}{
  \begin{question}{Q117}
    ข้อใดเป็นส่วนหนึ่งของ Vision & Goals
    \begin{choices}
        \wrongchoice{การใช้พลังงานหมุนเวียน}
        \correctchoice{การกำหนดเขต Smart City และเป้าหมายของเมือง}
        \wrongchoice{การทำระบบขนส่งอัจฉริยะ}
        \wrongchoice{การเก็บภาษีในเมือง}
    \end{choices}
  \end{question}
}

\element{smartcity}{
  \begin{question}{Q118}
    Infrastructure Plan เกี่ยวข้องกับเรื่องใด
    \begin{choices}
        \wrongchoice{การทำระบบฐานข้อมูลประชากร}
        \correctchoice{แผนการพัฒนาโครงสร้างพื้นฐาน ทั้งดิจิทัลและกายภาพ เช่น ขนส่ง พลังงาน สาธารณูปโภค}
        \wrongchoice{การจัดทำงบประมาณประจำปี}
        \wrongchoice{การส่งเสริมการท่องเที่ยว}
    \end{choices}
  \end{question}
}

\element{smartcity}{
  \begin{question}{Q119}
    City Data & Security ต้องมีอย่างน้อย 3 เรื่อง อะไรบ้าง
    \begin{choices}
        \correctchoice{City Data Catalogue, City Data Exchange, City Governance}
        \wrongchoice{Smart Economy, Smart Mobility, Smart Energy}
        \wrongchoice{E-Government, E-Banking, E-Commerce}
        \wrongchoice{Big Data, Blockchain, AI}
    \end{choices}
  \end{question}
}

\element{smartcity}{
  \begin{question}{Q120}
    คำว่า Smart City ตามนิยามของ UK Parliament หมายถึงอะไร
    \begin{choices}
        \wrongchoice{เมืองที่มีตึกสูงและระบบขนส่งสะดวก}
        \correctchoice{เมืองที่ใช้เทคโนโลยีและข้อมูลมาช่วยพัฒนาเมือง}
        \wrongchoice{เมืองที่มีประชากรหนาแน่นและเศรษฐกิจดี}
        \wrongchoice{เมืองที่มีมรดกทางวัฒนธรรมเยอะ}
    \end{choices}
  \end{question}
}

\element{smartcity}{
  \begin{question}{Q121}
    Smart Infrastructure มีบทบาทอย่างไรในเมืองอัจฉริยะ
    \begin{choices}
        \correctchoice{ระบบโครงสร้างพื้นฐานที่ช่วยสนับสนุนบริการดิจิทัล}
        \wrongchoice{การเลือกตั้งผู้ว่าราชการ}
        \wrongchoice{การใช้สมาร์ทโฟน}
        \wrongchoice{การจัดงานเทศกาล}
    \end{choices}
  \end{question}
}

\element{smartcity}{
  \begin{question}{Q122}
    ตัวอย่างของ Semiintelligent infrastructure คืออะไร
    \begin{choices}
        \wrongchoice{ระบบจราจรที่แนะนำเส้นทางให้ผู้ขับขี่}
        \correctchoice{แผนที่ที่บันทึกข้อมูลมลพิษหรือปริมาณรถในเมือง}
        \wrongchoice{ระบบไฟถนนที่ปรับแสงอัตโนมัติ}
        \wrongchoice{ระบบอาคารอัจฉริยะที่ควบคุมพลังงานเอง}
    \end{choices}
  \end{question}
}

\element{smartcity}{
  \begin{question}{Q123}
    ตัวอย่างของ Intelligent infrastructure คืออะไร
    \begin{choices}
        \wrongchoice{ระบบไฟถนนที่เปิดปิดเองตามเวลา}
        \correctchoice{ระบบขนส่งที่ตรวจจับการจราจรติดขัดและแจ้งผู้ขับ}
        \wrongchoice{แผนที่ที่บอกเฉพาะพื้นที่สีเขียว}
        \wrongchoice{ระบบประปาที่เก็บข้อมูลการใช้น้ำแต่ไม่ปรับเปลี่ยนอะไร}
    \end{choices}
  \end{question}
}

\element{smartcity}{
  \begin{question}{Q124}
    ตามแนวคิด Smart City ของไทย มีกี่ด้าน
    \begin{choices}
        \wrongchoice{5 ด้าน}
        \wrongchoice{6 ด้าน}
        \correctchoice{7 ด้าน}
        \wrongchoice{8 ด้าน}
    \end{choices}
  \end{question}
}

\element{smartcity}{
  \begin{question}{Q125}
    ใน 7 Dimensions of Smart City Solutions สิ่งที่บังคับว่าต้องมีคืออะไร
    \begin{choices}
        \wrongchoice{Smart Economy}
        \correctchoice{Smart Environment}
        \wrongchoice{Smart People}
        \wrongchoice{Smart Governance}
    \end{choices}
  \end{question}
}
